\section{Introduction}

The development of technology has made gadgets increasingly integrated
into learning activities. In Informatics classes, gadgets can be used to
access learning materials, educational applications, communication
platforms, information, and various interactive learning activities.
However, using a device more frequently in the classroom does not
automatically mean that learning becomes more effective. Therefore, data
is needed to examine how gadget use is related to students' learning
motivation and social interaction.

This article presents a learning journey from Introduction to Statistics,
Descriptive Statistics, Bivariate Analysis, to Basic Probability. The data
are based on a quantitative correlational study examining the use of
gadgets as a learning medium among Grade 12 students in Informatics II and
III at SMAK Yos Sudarso Batam during the 2024/2025 academic year.

The study involved \resN{} students selected using simple random sampling
\parencite{sugiyono2017}. Data were collected through a closed
questionnaire using a Likert scale and originally analyzed using IBM
SPSS 26.0 \parencite{ibm2021}; the analysis was independently reproduced
with the Python pipeline in this repository.
